\scanpage{8}
\uline{Morphismes radiciels} \quad Soit $f:X\mathrel{-}Y$ un morphisme de schémas. Conditions équivalentes

a) $f$ univ. \uline{injectif}

b) $\forall y\in Y$, $X_y$ est radiciel sur $k(y)$, i.e. \deleted{\unsure{est}} a au plus un point $x$, et $k(x)$ est une ext. radicielle de $k(y)$

\unsure{c|b}) $\forall$ corps $K$, $X(K)\longrightarrow Y(K)$ est \uline{injective}

d) $\delta_f:X\longrightarrow X\times_Y X$ est surjective

Les morphismes sont dits \uline{radiciels}. (Un \uline{monom.} est radiciel, donc une imm, telle $\delta_f$, l'est) Satisfont à b) et c] ($gf$ radiciel $\Rightarrow f$ radiciel)

\uline{Prop} Soit $f:X\mathrel{-}Y$ un morphisme entier radiciel \inserted{surjectif}. Alors c'est un \uline{homéom.} \deleted{\unsure{univ}} \uline{universel}.

On peut prouver la réciproque -- mais c'est nettement plus délicat

\uline{Exemples} 1) $k'/k$ ext. radicielle (p. ex. clôture parfaite)

$X_{k'}\mathrel{-}X_k$ est entier radiciel surjectif, donc un homéom (universel)

2) Courbe à cusp, et sa normalisée.

\uline{Morphismes univ. ouverts}

\uline{Condition néc.} Pour qu'un morphisme soit ouvert : il \uline{est générisant}

Cette condition est \uline{suffisante} si $f$ loc. de prés. finie

Donc morphismes univ. ouverts \deleted{sont} \uline{univ. générisants}, et réciproqt si $f$ est loc. de prés finie

Morphismes univ. ouverts et univ. générisants satisfont 1), 2).

% archive sheet number: 7
